% --- Archon physics formalization source begin ---
% archon:physics
% archon:covers PhyXMiniProblems/problem_phyx_mini_0708.lean
% archon:source-report reports/phyx_mini/problem_phyx_mini_0708.source.json

\chapter{Physics problem phyx\_mini\_0708}
\label{ch:PhyXMiniProblems_problem_phyx_mini_0708}

\paragraph{Problem source.}
\#\# Physical scenario

Suppose the earth suddenly came to a halt and ceased revolving around the sun. The gravitational force would then pull it directly into the sun.

\#\# Question

What would be the earth's speed as it crashed?

\#\# Answer choices

- A: A: \$5.13\textbackslash{}times10\textasciicircum{}5\textbackslash{}

- B: B: \$6.13\textbackslash{}times10\textasciicircum{}4\textbackslash{}

- C: C: \$6.13\textbackslash{}times10\textasciicircum{}5\textbackslash{}

- D: D: \$7.13\textbackslash{}times10\textasciicircum{}5\textbackslash{}

\#\# Figure caption (auxiliary; use the image as primary evidence)

The image is a diagram showing a "Before" and "After" scenario involving the Earth and the Sun, with measurements and annotations.

\#\#\# Before:
- The Earth is represented as a circle labeled "Earth."
- An arrow labeled \textbackslash{}( R\_e \textbackslash{}) points outward from the surface, indicating the Earth's radius.
- The velocity is marked as \textbackslash{}( v\_1 = 0 \textbackslash{}, \textbackslash{}text\{m/s\} \textbackslash{}).

\#\#\# After:
- The Sun is depicted as a larger yellow circle labeled "Sun."
- An arrow labeled \textbackslash{}( R\_s \textbackslash{}) emanates from the Sun, indicating its radius.
- A smaller circle representing the Earth is near the Sun, with a green arrow labeled \textbackslash{}( v\_2 \textbackslash{}) pointing towards the right, showing the new velocity direction.

\#\#\# Distances:
- The distance from the "Before" Earth's point to the Sun's surface in the "After" state is given as \textbackslash{}( r\_1 = 1.50 \textbackslash{}times 10\textasciicircum{}\{11\} \textbackslash{}, \textbackslash{}text\{m\} \textbackslash{}).
- The distance \textbackslash{}( r\_2 \textbackslash{}) from Earth's surface to Sun's surface after is calculated as \textbackslash{}( r\_2 = R\_s + R\_e = 7.02 \textbackslash{}times 10\textasciicircum{}8 \textbackslash{}, \textbackslash{}text\{m\} \textbackslash{}).

This image is likely illustrating concepts related to velocity change and distance in a celestial mechanics context.

\#\# Dataset metadata

Work and Energy; Physical Model Grounding Reasoning, Multi-Formula Reasoning

\paragraph{Recorded answer/context.}
C: C: \$6.13\textbackslash{}times10\textasciicircum{}5\textbackslash{}

\paragraph{Figure/image path.}
/root/proposal\_for\_physic/science-mango/phyx\_mini\_run/phyx\_data/test\_image/708.png

\paragraph{Formalization target.}
create a compiling Lean file with sorry bodies at `PhyXMiniProblems/problem\_phyx\_mini\_0708.lean`. The Lean declarations must preserve the physical quantities, dimensions or dimensional roles, figure labels, governing-law hypotheses, and final relation expressed by this problem.
Use Mathlib/Physlib names found through LeanExplore where available. If a physics API is missing, introduce faithful local abstractions rather than scalar placeholder aliases.

\begin{theorem}[Physics formalization target]
\label{thm:physics:phyx_mini_0708:target}
\lean{PhyXMiniProblems.ProblemPhyXMini0708.earthImpactSpeed_is_recordedChoiceC}
\uses{def:physics:phyx-mini-0708:phyxminiproblems-problemphyxmini0708-lengthinmeters, def:physics:phyx-mini-0708:phyxminiproblems-problemphyxmini0708-massinkilograms, def:physics:phyx-mini-0708:phyxminiproblems-problemphyxmini0708-speedinmeterspersecond, def:physics:phyx-mini-0708:phyxminiproblems-problemphyxmini0708-gravitationalconstantinsi, def:physics:phyx-mini-0708:phyxminiproblems-problemphyxmini0708-earthsunradialfallsetup, def:physics:phyx-mini-0708:phyxminiproblems-problemphyxmini0708-matchesstoppedearthscenario, def:physics:phyx-mini-0708:phyxminiproblems-problemphyxmini0708-matchesprimaryearthsunfigure, def:physics:phyx-mini-0708:phyxminiproblems-problemphyxmini0708-usestextbooksolarparameters, def:physics:phyx-mini-0708:phyxminiproblems-problemphyxmini0708-hasphysicalearthsunparameters, def:physics:phyx-mini-0708:phyxminiproblems-problemphyxmini0708-satisfiesnewtonianradialfallenergylaws, def:physics:phyx-mini-0708:phyxminiproblems-problemphyxmini0708-answerspeedinmeterspersecond, def:physics:phyx-mini-0708:phyxminiproblems-problemphyxmini0708-isuniqueclosestimpactspeedchoice}
The assigned autoformalize agent should translate this physics problem into Lean declarations in the covered file, with theorem and lemma proof bodies written as `by sorry`.
\end{theorem}
\begin{proof}
This is an autoformalization task, not a proof task. Produce statements that can later be proved by the physics prover without weakening the physical model.
\end{proof}

% --- Archon declaration topology begin ---
\section{Declaration topology}
\begin{definition}[Length Quantity]
  \label{def:physics:phyx-mini-0708:phyxminiproblems-problemphyxmini0708-lengthquantity}
  \lean{PhyXMiniProblems.ProblemPhyXMini0708.LengthQuantity}
  A nonnegative physical length, independent of a choice of units
\end{definition}
\begin{proof}
  This is the declared model definition of Length Quantity.
\end{proof}
\begin{definition}[Mass Quantity]
  \label{def:physics:phyx-mini-0708:phyxminiproblems-problemphyxmini0708-massquantity}
  \lean{PhyXMiniProblems.ProblemPhyXMini0708.MassQuantity}
  A nonnegative physical mass, independent of a choice of units
\end{definition}
\begin{proof}
  This is the declared model definition of Mass Quantity.
\end{proof}
\begin{definition}[Gravitational Constant Quantity]
  \label{def:physics:phyx-mini-0708:phyxminiproblems-problemphyxmini0708-gravitationalconstantquantity}
  \lean{PhyXMiniProblems.ProblemPhyXMini0708.GravitationalConstantQuantity}
  The Newtonian gravitational constant, with dimension length³ * mass⁻¹ * time⁻²
\end{definition}
\begin{proof}
  This is the declared model definition of Gravitational Constant Quantity.
\end{proof}
\begin{definition}[length In Meters]
  \label{def:physics:phyx-mini-0708:phyxminiproblems-problemphyxmini0708-lengthinmeters}
  \lean{PhyXMiniProblems.ProblemPhyXMini0708.lengthInMeters}
  \uses{def:physics:phyx-mini-0708:phyxminiproblems-problemphyxmini0708-lengthquantity}
  Read a physical length in SI metres
\end{definition}
\begin{proof}
  This is the declared model definition of length In Meters.
\end{proof}
\begin{definition}[mass In Kilograms]
  \label{def:physics:phyx-mini-0708:phyxminiproblems-problemphyxmini0708-massinkilograms}
  \lean{PhyXMiniProblems.ProblemPhyXMini0708.massInKilograms}
  \uses{def:physics:phyx-mini-0708:phyxminiproblems-problemphyxmini0708-massquantity}
  Read a physical mass in SI kilograms
\end{definition}
\begin{proof}
  This is the declared model definition of mass In Kilograms.
\end{proof}
\begin{definition}[speed In Meters Per Second]
  \label{def:physics:phyx-mini-0708:phyxminiproblems-problemphyxmini0708-speedinmeterspersecond}
  \lean{PhyXMiniProblems.ProblemPhyXMini0708.speedInMetersPerSecond}
  Read a nonnegative physical speed in SI metres per second
\end{definition}
\begin{proof}
  This is the declared model definition of speed In Meters Per Second.
\end{proof}
\begin{definition}[energy In Joules]
  \label{def:physics:phyx-mini-0708:phyxminiproblems-problemphyxmini0708-energyinjoules}
  \lean{PhyXMiniProblems.ProblemPhyXMini0708.energyInJoules}
  Read a physical energy in SI joules
\end{definition}
\begin{proof}
  This is the declared model definition of energy In Joules.
\end{proof}
\begin{definition}[gravitational Constant In SI]
  \label{def:physics:phyx-mini-0708:phyxminiproblems-problemphyxmini0708-gravitationalconstantinsi}
  \lean{PhyXMiniProblems.ProblemPhyXMini0708.gravitationalConstantInSI}
  \uses{def:physics:phyx-mini-0708:phyxminiproblems-problemphyxmini0708-gravitationalconstantquantity}
  Read the gravitational constant in m³ kg⁻¹ s⁻²
\end{definition}
\begin{proof}
  This is the declared model definition of gravitational Constant In SI.
\end{proof}
\begin{definition}[Celestial Body]
  \label{def:physics:phyx-mini-0708:phyxminiproblems-problemphyxmini0708-celestialbody}
  \lean{PhyXMiniProblems.ProblemPhyXMini0708.CelestialBody}
  The two physical bodies whose masses and radii enter the model
\end{definition}
\begin{proof}
  This is the declared model definition of Celestial Body.
\end{proof}
\begin{definition}[Fall State]
  \label{def:physics:phyx-mini-0708:phyxminiproblems-problemphyxmini0708-fallstate}
  \lean{PhyXMiniProblems.ProblemPhyXMini0708.FallState}
  The two states compared by conservation of mechanical energy
\end{definition}
\begin{proof}
  This is the declared model definition of Fall State.
\end{proof}
\begin{definition}[Radial Motion]
  \label{def:physics:phyx-mini-0708:phyxminiproblems-problemphyxmini0708-radialmotion}
  \lean{PhyXMiniProblems.ProblemPhyXMini0708.RadialMotion}
  Radial motion of the Earth's center relative to the Sun's center
\end{definition}
\begin{proof}
  This is the declared model definition of Radial Motion.
\end{proof}
\begin{definition}[Trajectory Geometry]
  \label{def:physics:phyx-mini-0708:phyxminiproblems-problemphyxmini0708-trajectorygeometry}
  \lean{PhyXMiniProblems.ProblemPhyXMini0708.TrajectoryGeometry}
  The idealization of the trajectory after orbital motion suddenly stops
\end{definition}
\begin{proof}
  This is the declared model definition of Trajectory Geometry.
\end{proof}
\begin{definition}[Interaction Model]
  \label{def:physics:phyx-mini-0708:phyxminiproblems-problemphyxmini0708-interactionmodel}
  \lean{PhyXMiniProblems.ProblemPhyXMini0708.InteractionModel}
  Which interactions are retained in the fall model
\end{definition}
\begin{proof}
  This is the declared model definition of Interaction Model.
\end{proof}
\begin{definition}[Central Body Approximation]
  \label{def:physics:phyx-mini-0708:phyxminiproblems-problemphyxmini0708-centralbodyapproximation}
  \lean{PhyXMiniProblems.ProblemPhyXMini0708.CentralBodyApproximation}
  Whether recoil of the much heavier central body is neglected
\end{definition}
\begin{proof}
  This is the declared model definition of Central Body Approximation.
\end{proof}
\begin{definition}[Horizontal Direction]
  \label{def:physics:phyx-mini-0708:phyxminiproblems-problemphyxmini0708-horizontaldirection}
  \lean{PhyXMiniProblems.ProblemPhyXMini0708.HorizontalDirection}
  Horizontal direction of the green impact-velocity arrow in the bitmap
\end{definition}
\begin{proof}
  This is the declared model definition of Horizontal Direction.
\end{proof}
\begin{definition}[Figure Label]
  \label{def:physics:phyx-mini-0708:phyxminiproblems-problemphyxmini0708-figurelabel}
  \lean{PhyXMiniProblems.ProblemPhyXMini0708.FigureLabel}
  Literal labels visible in the supplied primary image
\end{definition}
\begin{proof}
  This is the declared model definition of Figure Label.
\end{proof}
\begin{definition}[Earth Sun Fall Figure]
  \label{def:physics:phyx-mini-0708:phyxminiproblems-problemphyxmini0708-earthsunfallfigure}
  \lean{PhyXMiniProblems.ProblemPhyXMini0708.EarthSunFallFigure}
  \uses{def:physics:phyx-mini-0708:phyxminiproblems-problemphyxmini0708-celestialbody, def:physics:phyx-mini-0708:phyxminiproblems-problemphyxmini0708-fallstate, def:physics:phyx-mini-0708:phyxminiproblems-problemphyxmini0708-horizontaldirection, def:physics:phyx-mini-0708:phyxminiproblems-problemphyxmini0708-figurelabel}
  Qualitative evidence transcribed from image 708.png. The numerical values attached to r₁, r₂, and v₁ are related to physical quantities only by MatchesPrimaryEarthSunFigure below
\end{definition}
\begin{proof}
  This is the declared model definition of Earth Sun Fall Figure.
\end{proof}
\begin{definition}[Earth Sun Radial Fall Setup]
  \label{def:physics:phyx-mini-0708:phyxminiproblems-problemphyxmini0708-earthsunradialfallsetup}
  \lean{PhyXMiniProblems.ProblemPhyXMini0708.EarthSunRadialFallSetup}
  \uses{def:physics:phyx-mini-0708:phyxminiproblems-problemphyxmini0708-lengthquantity, def:physics:phyx-mini-0708:phyxminiproblems-problemphyxmini0708-massquantity, def:physics:phyx-mini-0708:phyxminiproblems-problemphyxmini0708-gravitationalconstantquantity, def:physics:phyx-mini-0708:phyxminiproblems-problemphyxmini0708-celestialbody, def:physics:phyx-mini-0708:phyxminiproblems-problemphyxmini0708-fallstate, def:physics:phyx-mini-0708:phyxminiproblems-problemphyxmini0708-radialmotion, def:physics:phyx-mini-0708:phyxminiproblems-problemphyxmini0708-trajectorygeometry, def:physics:phyx-mini-0708:phyxminiproblems-problemphyxmini0708-interactionmodel, def:physics:phyx-mini-0708:phyxminiproblems-problemphyxmini0708-centralbodyapproximation, def:physics:phyx-mini-0708:phyxminiproblems-problemphyxmini0708-earthsunfallfigure}
  The fields are independent physical quantities. In particular, earthSpeed at contact is not defined from an answer choice or from the desired formula
\end{definition}
\begin{proof}
  This is the declared model definition of Earth Sun Radial Fall Setup.
\end{proof}
\begin{definition}[Matches Stopped Earth Scenario]
  \label{def:physics:phyx-mini-0708:phyxminiproblems-problemphyxmini0708-matchesstoppedearthscenario}
  \lean{PhyXMiniProblems.ProblemPhyXMini0708.MatchesStoppedEarthScenario}
  \uses{def:physics:phyx-mini-0708:phyxminiproblems-problemphyxmini0708-speedinmeterspersecond, def:physics:phyx-mini-0708:phyxminiproblems-problemphyxmini0708-earthsunradialfallsetup}
  The prose scenario: the Earth has stopped revolving, is initially stationary, and subsequently moves radially toward a Sun treated as the fixed central body. No impact-speed value occurs here
\end{definition}
\begin{proof}
  This is the declared model definition of Matches Stopped Earth Scenario.
\end{proof}
\begin{definition}[Matches Primary Earth Sun Figure]
  \label{def:physics:phyx-mini-0708:phyxminiproblems-problemphyxmini0708-matchesprimaryearthsunfigure}
  \lean{PhyXMiniProblems.ProblemPhyXMini0708.MatchesPrimaryEarthSunFigure}
  \uses{def:physics:phyx-mini-0708:phyxminiproblems-problemphyxmini0708-lengthinmeters, def:physics:phyx-mini-0708:phyxminiproblems-problemphyxmini0708-earthsunradialfallsetup}
  Facts read from the primary bitmap. The two distances are center-to-center separations; at contact this equals the sum of the two physical radii. These are input geometry and measurement readouts, not the requested impact speed
\end{definition}
\begin{proof}
  This is the declared model definition of Matches Primary Earth Sun Figure.
\end{proof}
\begin{definition}[Uses Textbook Solar Parameters]
  \label{def:physics:phyx-mini-0708:phyxminiproblems-problemphyxmini0708-usestextbooksolarparameters}
  \lean{PhyXMiniProblems.ProblemPhyXMini0708.UsesTextbookSolarParameters}
  \uses{def:physics:phyx-mini-0708:phyxminiproblems-problemphyxmini0708-massinkilograms, def:physics:phyx-mini-0708:phyxminiproblems-problemphyxmini0708-gravitationalconstantinsi, def:physics:phyx-mini-0708:phyxminiproblems-problemphyxmini0708-earthsunradialfallsetup}
  Independent textbook constants required to evaluate the numerical answer. The source image supplies the geometry but does not print G or the solar mass, so their use is made explicit rather than hidden in a local definition
\end{definition}
\begin{proof}
  This is the declared model definition of Uses Textbook Solar Parameters.
\end{proof}
\begin{definition}[Has Physical Earth Sun Parameters]
  \label{def:physics:phyx-mini-0708:phyxminiproblems-problemphyxmini0708-hasphysicalearthsunparameters}
  \lean{PhyXMiniProblems.ProblemPhyXMini0708.HasPhysicalEarthSunParameters}
  \uses{def:physics:phyx-mini-0708:phyxminiproblems-problemphyxmini0708-lengthinmeters, def:physics:phyx-mini-0708:phyxminiproblems-problemphyxmini0708-massinkilograms, def:physics:phyx-mini-0708:phyxminiproblems-problemphyxmini0708-gravitationalconstantinsi, def:physics:phyx-mini-0708:phyxminiproblems-problemphyxmini0708-earthsunradialfallsetup}
  Physical-domain conditions needed to use and cancel the model quantities
\end{definition}
\begin{proof}
  This is the declared model definition of Has Physical Earth Sun Parameters.
\end{proof}
\begin{definition}[Satisfies Newtonian Radial Fall Energy Laws]
  \label{def:physics:phyx-mini-0708:phyxminiproblems-problemphyxmini0708-satisfiesnewtonianradialfallenergylaws}
  \lean{PhyXMiniProblems.ProblemPhyXMini0708.SatisfiesNewtonianRadialFallEnergyLaws}
  \uses{def:physics:phyx-mini-0708:phyxminiproblems-problemphyxmini0708-lengthinmeters, def:physics:phyx-mini-0708:phyxminiproblems-problemphyxmini0708-massinkilograms, def:physics:phyx-mini-0708:phyxminiproblems-problemphyxmini0708-speedinmeterspersecond, def:physics:phyx-mini-0708:phyxminiproblems-problemphyxmini0708-energyinjoules, def:physics:phyx-mini-0708:phyxminiproblems-problemphyxmini0708-gravitationalconstantinsi, def:physics:phyx-mini-0708:phyxminiproblems-problemphyxmini0708-earthsunradialfallsetup}
  The kinetic-energy formula, point-mass central gravitational potential, and conservation of mechanical energy. These are governing laws at two independent states; they do not contain the solved speed-squared relation or a numerical impact speed
\end{definition}
\begin{proof}
  This is the declared model definition of Satisfies Newtonian Radial Fall Energy Laws.
\end{proof}
\begin{definition}[Answer Choice]
  \label{def:physics:phyx-mini-0708:phyxminiproblems-problemphyxmini0708-answerchoice}
  \lean{PhyXMiniProblems.ProblemPhyXMini0708.AnswerChoice}
  Labels of the four impact-speed choices in the source
\end{definition}
\begin{proof}
  This is the declared model definition of Answer Choice.
\end{proof}
\begin{definition}[answer Speed In Meters Per Second]
  \label{def:physics:phyx-mini-0708:phyxminiproblems-problemphyxmini0708-answerspeedinmeterspersecond}
  \lean{PhyXMiniProblems.ProblemPhyXMini0708.answerSpeedInMetersPerSecond}
  \uses{def:physics:phyx-mini-0708:phyxminiproblems-problemphyxmini0708-answerchoice}
  Metres-per-second readout printed beside each answer label
\end{definition}
\begin{proof}
  This is the declared model definition of answer Speed In Meters Per Second.
\end{proof}
\begin{definition}[Is Unique Closest Impact Speed Choice]
  \label{def:physics:phyx-mini-0708:phyxminiproblems-problemphyxmini0708-isuniqueclosestimpactspeedchoice}
  \lean{PhyXMiniProblems.ProblemPhyXMini0708.IsUniqueClosestImpactSpeedChoice}
  \uses{def:physics:phyx-mini-0708:phyxminiproblems-problemphyxmini0708-speedinmeterspersecond, def:physics:phyx-mini-0708:phyxminiproblems-problemphyxmini0708-earthsunradialfallsetup, def:physics:phyx-mini-0708:phyxminiproblems-problemphyxmini0708-answerchoice, def:physics:phyx-mini-0708:phyxminiproblems-problemphyxmini0708-answerspeedinmeterspersecond}
  A choice is uniquely closest when every distinct printed alternative has strictly larger absolute error from the independent physical impact speed
\end{definition}
\begin{proof}
  This is the declared model definition of Is Unique Closest Impact Speed Choice.
\end{proof}
\begin{lemma}[speed Squared Change eq newtonian Potential Drop]
  \label{lem:physics:phyx-mini-0708:phyxminiproblems-problemphyxmini0708-speedsquaredchange-eq-newtonianpotentialdrop}
  \lean{PhyXMiniProblems.ProblemPhyXMini0708.speedSquaredChange_eq_newtonianPotentialDrop}
  \uses{def:physics:phyx-mini-0708:phyxminiproblems-problemphyxmini0708-lengthinmeters, def:physics:phyx-mini-0708:phyxminiproblems-problemphyxmini0708-massinkilograms, def:physics:phyx-mini-0708:phyxminiproblems-problemphyxmini0708-speedinmeterspersecond, def:physics:phyx-mini-0708:phyxminiproblems-problemphyxmini0708-gravitationalconstantinsi, def:physics:phyx-mini-0708:phyxminiproblems-problemphyxmini0708-earthsunradialfallsetup, def:physics:phyx-mini-0708:phyxminiproblems-problemphyxmini0708-hasphysicalearthsunparameters, def:physics:phyx-mini-0708:phyxminiproblems-problemphyxmini0708-satisfiesnewtonianradialfallenergylaws}
  Energy conservation gives the general change in squared speed between the initial and contact separations
\end{lemma}
\begin{proof}
  The conclusion follows from the governing relations and figure readouts named in the statement of speed Squared Change eq newtonian Potential Drop.
\end{proof}
\begin{lemma}[impact Speed Squared eq newtonian Potential Drop]
  \label{lem:physics:phyx-mini-0708:phyxminiproblems-problemphyxmini0708-impactspeedsquared-eq-newtonianpotentialdrop}
  \lean{PhyXMiniProblems.ProblemPhyXMini0708.impactSpeedSquared_eq_newtonianPotentialDrop}
  \uses{def:physics:phyx-mini-0708:phyxminiproblems-problemphyxmini0708-lengthinmeters, def:physics:phyx-mini-0708:phyxminiproblems-problemphyxmini0708-massinkilograms, def:physics:phyx-mini-0708:phyxminiproblems-problemphyxmini0708-speedinmeterspersecond, def:physics:phyx-mini-0708:phyxminiproblems-problemphyxmini0708-gravitationalconstantinsi, def:physics:phyx-mini-0708:phyxminiproblems-problemphyxmini0708-earthsunradialfallsetup, def:physics:phyx-mini-0708:phyxminiproblems-problemphyxmini0708-matchesstoppedearthscenario, def:physics:phyx-mini-0708:phyxminiproblems-problemphyxmini0708-hasphysicalearthsunparameters, def:physics:phyx-mini-0708:phyxminiproblems-problemphyxmini0708-satisfiesnewtonianradialfallenergylaws}
  With the initial speed zero, the general energy relation is the impact law
\end{lemma}
\begin{proof}
  The conclusion follows from the governing relations and figure readouts named in the statement of impact Speed Squared eq newtonian Potential Drop.
\end{proof}
% --- Archon declaration topology end ---
% --- Archon physics formalization source end ---
